\section{Full compact spatial realization through the third transition}
\label{tts:section}

We continue the actual third-growth fields through the entire third
transition just proved. The second vorticity now revives with negative
sign. Both older velocities and temperatures enter the spatial
equations. The construction below uses the actual coefficient solution
and proves its complete spatial realization, including a smooth finite
extension of its forces. Its endpoint is the third transition endpoint;
no third return is asserted here.

\subsection{The original objects and a uniform material bound}

The physical equations, coordinates, and orientation remain
\begin{align}\label{tts:PDE}
 \theta_t+u\cdot\nabla\theta&=d\Delta\theta+f_\theta,
 &x&\in\mathbb R^2,\nonumber\\
 u_t+(u\cdot\nabla)u&=d\Delta u+\theta e_2+f_u,
 &\operatorname{div}u&=0,\qquad p=0,\\
 J&=\begin{pmatrix}0&-1\\1&0\end{pmatrix},
 & e(s)&=(\sin s,\cos s).\nonumber
\end{align}
Here $d>0$ is the common physical viscosity and thermal diffusivity;
$A_0\geq1$, $\nu=\sqrt{A_0}$, and
$\mu=\lambda_0=\lceil\nu/2\rceil$ retain their original values.
In particular $\nu$ is a time scale. Retain
$K_j=\mu\widehat\lambda_j$, the exact source integer recursion,
the source step $\eta$, and the actual scales $\sigma_1,\sigma_2$.
Put
\begin{equation}\label{tts:times}
 t_b=t_2^b,\quad t_a=t_3^a,\quad
 s=s_3=L_3\sigma_1/\sigma_2,\quad
 \Gamma=\Gamma_3=\sigma_2\sin s,\quad
 t_1=t_a+\Gamma^{-1},\quad I=[t_b,t_1].
\end{equation}
These are the actual second return, third threshold, and third
transition endpoint already constructed. Their proved bounds include
$0<t_a-t_b<14/\sigma_1$, $L_3\geq16384$,
$0<s\leq a:=\sin s_2\leq1/512$, and
$\Gamma\geq(15/16)L_3\sigma_1$. Thus
\begin{equation}\label{tts:duration}
 |I|=(t_a-t_b)+\Gamma^{-1}<15/\sigma_1.
\end{equation}
All earlier activation weights are one on $I$; the third is
$w_3(t)=\eta(\Gamma(t-t_b))$ and is one before $t_a$.
The exact source control on $I$ is
$\dot\alpha=(1-\lambda)\mathcal F_2+\lambda\mathcal F_3$,
where $\lambda(t)=\eta(\Gamma(t-t_a))$ is zero during growth.
The two switches $w_3,\lambda$ have different original start times.

Use the determinant coordinates of the third-transition theorem:
\begin{equation}\label{tts:geometry}
 \begin{gathered}
 h_0=\cos s_2,\quad D=h^2+a^2,\quad
 r_b=|\zeta_2(t_b)|,\quad b=r_b\sin s,\quad
 R=|\zeta_3|^2=(g^2+b^2)/D,\\
 \zeta_1=\frac{(h\sin\vartheta-a\cos\vartheta,
                     h\cos\vartheta+a\sin\vartheta)}{\sqrt D},\\
 \zeta_2=\sqrt D(\sin\vartheta,\cos\vartheta),\\
 \zeta_3=\frac{(g\sin\vartheta+b\cos\vartheta,
                     g\cos\vartheta-b\sin\vartheta)}{\sqrt D},
 \qquad\vartheta=su.
 \end{gathered}
\end{equation}
The symbol $D$ in this display is a scalar, not a matrix.
The actual growth and transition estimates give, on all of $I$,
\begin{equation}\label{tts:box}
 \begin{gathered}
 h_0\leq h<13/4,\quad 1\leq D<11,\quad 3/4<g<5,
 \quad e^{-1/64}\leq R<3,\\
 0\leq u\leq640a/L_3<1/1024,\quad
 \cos\vartheta>15/16,\quad
 g\cos\vartheta-b\sin\vartheta>2/3,\\
 0<\Omega_1<10\sigma_1,\quad
 -7040a\sigma_2/L_3\leq\Omega_2\leq0,\quad
 0<-\Theta_1\leq2\sigma_1^2/K_1,\\
 0<-\Theta_2\leq3\sigma_2^2/K_2,\quad
 0<-\Theta_3\leq e^4P_3^*,\quad
 |\Omega_3|\leq3e^4K_3P_3^*/\sigma_2,\\
 P_3^*=(A_0/\mu)\widehat\lambda_3^{-7/8},\qquad
 1\leq r_b\leq\sqrt{10}.
 \end{gathered}
\end{equation}
For the bound on $R$, growth has $1\leq R\leq2$, and the
transition changes $\log|\zeta_3|$ by less than $1/128$.
The amplitude bounds follow from the actual growth target, the
transition gain at most four, and the proved quotient at most three.
They are consequences of the constructed trajectory, not additional
conditions imposed on an unknown flow.

For completeness the exact material equations for all three layers are
\begin{align}\label{tts:material}
 G_{<j}&=-A_0R_\alpha e_2+\sum_{i<j}w_iK_i\Theta_i\zeta_i,
 &D_{<j}&=\dot\alpha J+
   \sum_{i<j}\frac{w_i\Omega_i}{r_i^2}J\zeta_i\otimes\zeta_i,
 \nonumber\\
 \dot M_j&=D_{<j}M_j,& M_j(t_{0j})&=I,
 \quad\zeta_j=M_j^{-T}p_j,\quad r_j=|\zeta_j|,\\
 \dot\Theta_j&=-\frac{J\zeta_j\cdot G_{<j}}{K_jr_j^2}\Omega_j
                       -dK_j^2r_j^2\Theta_j,
 &\dot\Omega_j&=K_j\zeta_{j,1}\Theta_j-dK_j^2r_j^2\Omega_j.
 \nonumber
\end{align}
Here $p_1=e(s_*)$, $p_2=e(s_2)$, $p_3=e(s)$, and
$t_{03}=t_b$; the first two original insertion times are unchanged.
The scalar $r_j^2$ is strictly positive. Each matrix $D_{<j}$ has
trace zero, since $\zeta_i\cdot J\zeta_i=0$.
The exact matrices are
\begin{equation}\label{tts:matrices}
 \begin{gathered}
 M_1=R_\alpha,\qquad
 M_2=R_{\alpha-s_*}
       \begin{pmatrix}1&-(h-h_0)/a\\0&1\end{pmatrix},\\
 \mathcal L(t)=(\zeta_2(t)\ \zeta_3(t)),\qquad
 M_3=\mathcal L(t)^{-T}\mathcal L(t_b)^T,\quad
 M_3^{-1}=\mathcal L(t_b)^{-T}\mathcal L(t)^T.
 \end{gathered}
\end{equation}
Indeed $\zeta_2$ is annihilated by the transpose of its own
added shear. Thus both columns of $\mathcal L$ solve
$\dot z=-D_{<3}^Tz$. Their determinant is
$\det\mathcal L=-b$, including at birth. Differentiating the
last two expressions proves $\dot M_3=D_{<3}M_3$,
$M_3(t_b)=I$, and $M_3^{-T}p_3=\zeta_3$.
The same calculation, or the displayed triangular form, proves
the first two material equations and their original data.
All matrices have determinant one, and hence define global linear
diffeomorphisms of the original $\mathbb R^2$.

Let $\|\cdot\|$ be Euclidean operator norm and
$\|\cdot\|_F$ Frobenius norm. For a two by two matrix its
adjugate has the same Frobenius norm: it permutes its four entries
and changes two signs. Consequently
\[
 \|\mathcal L^{-1}\|\leq\|\mathcal L\|_F/b,\quad
 \|\mathcal L\|_F^2=D+R<14,\quad
 \|\mathcal L(t_b)\|_F^2=r_b^2+1\leq11.
\]
Applying these inequalities to both products in
\eqref{tts:matrices} gives the uniform bounds
\begin{equation}\label{tts:matrixbounds}
 \begin{gathered}
 N_1=1,\qquad N=1+(3-h_0)/a,\qquad
 N_2=\widetilde N=1+(13/4-h_0)/a\leq2N,\\
 N_3=\mathcal M=16/\sin s,\qquad
 \|M_j\|,\|M_j^{-1}\|\leq N_j\quad(1\leq j\leq3).
 \end{gathered}
\end{equation}
For $j=3$ the product bound is $\sqrt{154}/b<16/\sin s$.
For $j=2$, write its triangular factor as $I-qe_1\otimes e_2$
and its inverse as $I+qe_1\otimes e_2$, so both norms are at
most $1+|q|$. Finally
$2N-\widetilde N=1+(11/4-h_0)/a>0$.
The third estimate is valid when $\Omega_2$ is nonzero; it does
not use the shorter growth-only shear formula.

\subsection{The exact source supports remain nested}

The birth bound $P_{2,b}\leq e^{18}(A_0/\mu)
\widehat\lambda_2^{-7/8}$ gives
$A_{2,b}=K_2r_bP_{2,b}\leq\sqrt{10}e^{18}\nu^2
\widehat\lambda_2^{1/8}$. The already proved scale relation
$\sigma_2^2\leq5A_{2,b}/4$ therefore gives
\begin{equation}\label{tts:sourceconstants}
 \begin{gathered}
 C_2=e^9\sqrt{5\sqrt{10}/4},\qquad
 \sigma_2\leq C_2\nu\widehat\lambda_2^{1/16},\\
 \sigma_1\geq\sqrt e\,\nu\widehat\lambda_1^{1/16},\qquad
 C_M=\frac{256C_2}{15\sqrt e},\qquad
 \mathcal M\leq C_M\widehat\lambda_1^{(Q_2-1)/16}.
 \end{gathered}
\end{equation}
For the last estimate use
$\sin s\geq15L_3\sigma_1/(16\sigma_2)$, $L_3\geq1$,
and the unchanged recursion
$\widehat\lambda_2=\widehat\lambda_1^{Q_2}$.
With $N\leq C_N\widehat\lambda_1^{1/16}$, the strengthened
actual source selection includes the strict inequalities
\begin{align}\label{tts:sourceexponents}
 \widehat\lambda_1^{47Q_2/16-3}&>4C_NC_M,
 &\widehat\lambda_1^{(31Q_2+1)/16}&>\sqrt{11}C_M/s_0.
\end{align}
Their exponents are positive because $Q_2\geq202$.
Multiplication by the corresponding negative frequency powers gives
\begin{equation}\label{tts:nestingnumbers}
 \begin{gathered}
 2N_2N_3\widehat\lambda_2^{-3}
       <\widehat\lambda_1^{-3},\qquad
 \sqrt{11}N_3\widehat\lambda_2^{-2}<s_0,\\
 N_2\widehat\lambda_1^{-3}<s_0/(2C_\ell),\qquad
 N_2\widehat\lambda_1^{-2}<s_0.
 \end{gathered}
\end{equation}
For the first, the exponent of $\widehat\lambda_1$ in
$4C_NC_M\widehat\lambda_1^{Q_2/16-3Q_2+3}$ is
$-(47Q_2/16-3)$. For the second its exponent is
$(Q_2-1)/16-2Q_2=-(31Q_2+1)/16$.
The last two follow from the earlier actual first-source margins
$N\widehat\lambda_1^{-3}<s_0/(4C_\ell)$ and
$N\widehat\lambda_1^{-2}<s_0/2$, together with $N_2\leq2N$.
Thus the doubling needed in the third transition is absorbed by
proved strict margins.

Retain the original profile, primitive, and spatial cutoffs. Namely
$F(s)=\sin s+\chi_F(s)(s-\sin s)$ on $[-\pi,\pi]$ is
extended periodically; $\chi_F$ is even, smooth, supported away
from the endpoints, and equals one on $[-s_0,s_0]$, $0<s_0<1$.
The even periodic primitive is uniquely specified by $P'=F$ and
$\int_{-\pi}^{\pi}P=0$. On the affine cell it is
$P(s)=P(0)+s^2/2$, with its original constant $P(0)$ retained.
The radial $f\in C_c^\infty(B_1)$ equals one on $B_{1/2}$.
Keep the exact source radii and material envelopes
\begin{equation}\label{tts:envelopes}
 \begin{gathered}
 \ell_1=s_0/(C_\ell\mu),\quad
 \ell_2=\widehat\lambda_1^{-3}/\mu,\quad
 \ell_3=\widehat\lambda_2^{-3}/\mu,\quad C_\ell\geq4,\\
 g_j(x,t)=f(M_j^{-1}x/\ell_j),\qquad
 k_j=K_j\zeta_j,\quad \xi_j=k_j\cdot x,
 \quad\rho_j=|k_j|^2.
 \end{gathered}
\end{equation}
The vectors $k_j$ here are physical phases, not the source integer
seed exponents. Every third field and derivative is supported in
$|x|\leq N_3\ell_3$. At such points
\[
 |M_2^{-1}x|\leq N_2N_3\ell_3<\ell_2/2,\qquad
 |K_2\zeta_2\cdot x|\leq K_2\sqrt{11}N_3\ell_3<s_0.
\]
Thus the full second field is affine on an open neighborhood of
the third support. Every point of the second support satisfies
$|x|\leq N_2\ell_2<\ell_1/2$ and
$|K_1\zeta_1\cdot x|<s_0$ by the last two inequalities in
\eqref{tts:nestingnumbers}; $M_1$ is a rotation. The first
support lies in $B_{\ell_1}$, where the original radial base
cutoff is one and $|\mu R_\alpha e_2\cdot x|<s_0$.
This proves every inclusion in the original laboratory coordinates.
Strictness ensures an open affine neighborhood for every support
and every derivative support, not only an identity at its center.

\subsection{All fields, original initial data, and full forces}

Retain the original radial $\chi_0\in C_c^\infty(B_{R_0})$,
equal to one on $B_2$, and the radial
$H_0\in C_c^\infty(B_{R_H})$, $R_H>R_0+1$, with
$H_0=|x|^2/2$ on $B_{R_0+1}$. Define
\begin{align}\label{tts:base}
 \theta_{\rm in}&=-\frac{A_0}{\mu}\sin(\mu x_2)\chi_0,
 &\theta_{\rm ref}&=-\frac{A_0}{\mu}F(\mu x_2)\chi_0,
 \nonumber\\
 \theta_b&=-\frac{A_0}{\mu}F(\mu R_\alpha e_2\cdot x)\chi_0
       +(1-\eta(\nu t))(\theta_{\rm in}-\theta_{\rm ref}),
 &\psi_b&=\dot\alpha H_0,\quad u_b=J\nabla\psi_b.
\end{align}
The angle is zero before $1/\nu$. Thus this is exactly the
original sine initial datum and zero initial velocity. The
interpolation has finished before the first insertion and hence
before $I$. For $j=1,2,3$ define the active increments
\begin{equation}\label{tts:fields}
 \begin{gathered}
 T_j=w_j\Theta_j,\quad Q_j=w_j\Omega_j/\rho_j,\quad
 V_j=T_jF(\xi_j)g_j,\quad\Psi_j=Q_jP(\xi_j)g_j,\quad
 v_j=J\nabla\Psi_j,\\
 \theta=\theta_b+\sum_{j=1}^3V_j,\qquad
 \psi=\psi_b+\sum_{j=1}^3\Psi_j,\qquad u=J\nabla\psi.
 \end{gathered}
\end{equation}
For $t\leq t_{0j}$ the $j$th increment is zero. Flatness of the
original $\eta$ makes this zero extension smooth. Before $t_b$
these formulas give the earlier constructed spatial solution
exactly. On $I$ they use the actual revived $\Omega_2$ and its
corresponding streamfunction; none of the older fields is frozen.

For each layer suppress its index temporarily and set
$k=k_j$, $\rho=\rho_j$, $\xi=k\cdot x$, $g=g_j$,
$D_*=D_{<j}$, $G_*=G_{<j}$. All profiles in the next formulas
are evaluated at $\xi$. Set
\begin{equation}\label{tts:Qrest}
 Q_{\rm r}=\frac{\dot w\Omega+w k_1\Theta}{\rho}
                    -Q\frac{\dot\rho}{\rho},\qquad
 \dot\rho=-2k\cdot D_*^Tk,
 \qquad\dot Q=Q_{\rm r}-d\rho Q.
\end{equation}
The exact force increments, with no pressure reconstruction, are
\begin{align}\label{tts:scalarforce}
 E_{\theta,j}={}&\dot w\Theta Fg+QP(J\nabla g\cdot G_*)
           +QT(F^2-PF')g(Jk\cdot\nabla g)\nonumber\\
 &-dT\{\rho(F+F'')g+2F'k\cdot\nabla g+F\Delta g\},\\
 E_{u,j}={}&2D_*v+Q_{\rm r}J(kFg+P\nabla g)
          +(\nabla v)v-TFg e_2-d(\rho v+\Delta v).
 \label{tts:vectorforce}
\end{align}
The complete spatial differentiations in these expressions are
\begin{align}\label{tts:spatialjets}
 v={}&QJ(kFg+P\nabla g),\nonumber\\
 \nabla v={}&QJ\{k\otimes k F'g+
 F(k\otimes\nabla g+\nabla g\otimes k)+P\nabla^2g\},\\
 \Delta V={}&T\{\rho F''g+2F'k\cdot\nabla g+F\Delta g\},
 \nonumber\\
 \Delta v={}&QJ\{\rho kF''g+\rho F'\nabla g
       +2F'k(k\cdot\nabla g)+2F(\nabla^2g)k
       +Fk\Delta g+P\nabla\Delta g\}.
 \nonumber
\end{align}
Keep the full base residuals
\begin{align}\label{tts:baseforce}
 f_{\theta,b}&=\nu\eta'(\nu t)(\theta_{\rm ref}-\theta_{\rm in})
                                      -d\Delta\theta_b,\nonumber\\
 f_{u,b}&=\ddot\alpha J\nabla H_0
    +\dot\alpha^2\bigl(\nabla(J\nabla H_0)\bigr)J\nabla H_0
    -\theta_b e_2-d\dot\alpha J\nabla\Delta H_0,\\
 f_\theta&=f_{\theta,b}+\sum_{j=1}^3E_{\theta,j},\qquad
 f_u=f_{u,b}+\sum_{j=1}^3E_{u,j}.\nonumber
\end{align}

\begin{proof}[Verification of every equation and interaction]
On a $j$th support the proved nesting makes the full older fields
exactly $G_*\cdot x$ and $D_*x$. For
$\mathcal D_j=\partial_t+D_*x\cdot\nabla$, differentiating
$M_j^{-1}$ gives $\mathcal D_j(M_j^{-1}x)=0$.
Likewise $\dot k=-D_*^Tk$ implies $\mathcal D_j\xi=0$.
Therefore $\mathcal D_jg=0$, $\mathcal D_jV=\dot T Fg$,
and $\mathcal D_j\Psi=\dot QPg$. Direct contractions give
\[
 \begin{split}
 v\cdot G_*&=QFg(Jk\cdot G_*)+QP(J\nabla g\cdot G_*),\\
 v\cdot\nabla V&=QT(F^2-PF')g(Jk\cdot\nabla g).
 \end{split}
\]
The second identity uses $Jk\cdot k=0$,
$J\nabla g\cdot\nabla g=0$, and
$J\nabla g\cdot k=-Jk\cdot\nabla g$.
The coupling in $\dot\Theta$ in \eqref{tts:material}
cancels exactly the first term of the first identity. Adding
$\dot w\Theta Fg$ and subtracting the complete $d\Delta V$
proves \eqref{tts:scalarforce}.

For an arbitrary trace-free two by two $D_*$, direct entrywise
multiplication gives $D_*J+JD_*^T=0$. Hence
\[
 \mathcal D_j(J\nabla\Psi)=D_*v+J\nabla(\mathcal D_j\Psi).
\]
The reverse velocity interaction is
$(v\cdot\nabla)(D_*x)=D_*v$. Adding it, the exact
self-advection, and the negative buoyancy and diffusion terms
gives \eqref{tts:vectorforce}, including its factor two.
The product rule applied to $\Psi=QPg$ proves all entries of
\eqref{tts:spatialjets}; differentiating
$\Delta\Psi=Q(\rho F'g+2Fk\cdot\nabla g+P\Delta g)$
in particular produces all six vector Laplacian groups.

These are full ordered interactions. In the third increment,
$\mathcal D_3$ contains both $v_1\cdot\nabla V_3$ and
$v_2\cdot\nabla V_3$, and both corresponding advections of
$v_3$. The term $v_3\cdot G_{<3}$ contains both
$v_3\cdot\nabla V_1$ and $v_3\cdot\nabla V_2$; the
term $D_{<3}v_3$ contains both reverse velocity interactions.
The same statements for $j=2$ include all pairs of the first
two layers, and $j=1$ includes the base. Thus adding in order
$1,2,3$ counts every ordered pair exactly once. Outside a layer's
support its derivatives vanish; the identities consequently telescope
on all of $\mathbb R^2$, not merely on an affine neighborhood.

The base radial cutoff is transported by $\dot\alpha Jx$,
and so is the base phase $\mu R_\alpha e_2\cdot x$.
On its scalar support $u_b=\dot\alpha Jx$; before $1/\nu$
the angle and velocity are zero. This proves the scalar base
formula, including the initial interpolation. Differentiating
$u_b=\dot\alpha J\nabla H_0$ proves its vector formula
globally. The sum therefore solves \eqref{tts:PDE} on
$[0,t_1]$ with the original data. Since $u=J\nabla\psi$,
equality of mixed spatial partial derivatives gives divergence zero.
\end{proof}

The exact spatial coordinate correspondence is also retained.
For each $j$ put $C_j=M_j^{-1}M_j^{-T}$ and
$\widetilde g_j(y)=f(y/\ell_j)$. On a function of independent
variables $(\xi,y)$, pulled back by
$(\xi,y)=(K_j\zeta_j\cdot x,M_j^{-1}x)$, the chain rule gives
\begin{equation}\label{tts:metric}
 \begin{split}
 \nabla_x&=K_j\zeta_j\partial_\xi+M_j^{-T}\nabla_y,\\
 \Delta_x&=K_j^2r_j^2\partial_\xi^2
   +2K_j(C_jp_j\cdot\nabla_y)\partial_\xi
   +\sum_{l,m}(C_j)_{lm}\partial_{y_l}\partial_{y_m}.
 \end{split}
\end{equation}
Multiplying the commuting constant spatial derivatives proves the
second equality from the first; in particular it retains all mixed
terms and the actual time-dependent metric.

On an affine cell $g_j=1$, $|\xi_j|<s_0$, the exact fields are
\begin{equation}\label{tts:core}
 V_j=T_jK_j\zeta_j\cdot x,\qquad
 v_j=\frac{w_j\Omega_j}{r_j^2}J\zeta_j\otimes\zeta_j\,x,
 \qquad\Delta V_j=\Delta v_j=0.
\end{equation}
This follows by differentiating the original
$Q_j(P(0)+\xi_j^2/2)$. The constant remains in the full cutoff
field; it was not discarded from its definition. The damping forces
in the core remain exactly $-d\rho_jV_j$ and $-d\rho_jv_j$.
On an uncut sine-profile region instead, $F''=-F$ gives
$\rho_jV_j+\Delta V_j=0$ and
$\rho_jv_j+\Delta v_j=0$. All other vector residual terms in
\eqref{tts:vectorforce} remain. Thus the core and sine calculations
give the exact map from the damped amplitude trajectories to the
released profile, including the forces that the profile requires.

\subsection{Quantitative spatial costs for every layer and both forces}

For any scalar, vector or tensor field $U$, write
$|U|_m=\sup_x\|D_x^mU\|$ in the Euclidean multilinear norm.
Let $[W]_q=\sup_\xi|W^{(q)}(\xi)|$ for a periodic profile and
$f_m=\sup_y\|D^mf(y)\|$. Define fixed bounds on $I$ by
\begin{equation}\label{tts:normdata}
 \begin{gathered}
 r_{-,j}^2=(1,1,63/64)_j,\qquad
 r_{+,j}^2=(1,11,3)_j,\qquad \kappa_j=K_jr_{+,j},\\
 U_{\Theta,j}=(2\sigma_1^2/K_1,\ 3\sigma_2^2/K_2,
                                      \ e^4P_3^*)_j,\\
 U_{\Omega,j}=(10\sigma_1,\ 7040a\sigma_2/L_3,
                             \ 3e^4K_3P_3^*/\sigma_2)_j,\\
 W_{j,1}=(0,0,\Gamma[\eta']_0)_j,\qquad
 \vartheta_*=640as/L_3,\qquad
 A_*=270a^2\sigma_1+6aU_{\Omega,2}\vartheta_*.
 \end{gathered}
\end{equation}
The lower third length follows from
$e^{-1/64}\geq1-1/64=63/64$. To verify the control bound
$|\dot\alpha|\leq A_*$, write
\[
 y=\frac{ag+hb}{D},\quad
 \mathcal F_2=\frac{a\Omega_1(h\tan\vartheta-a)}D,\quad
 \mathcal F_3=
 \frac{\Omega_1y(h\sin\vartheta-a\cos\vartheta)
                  +b\Omega_2\sin\vartheta}
      {g\cos\vartheta-b\sin\vartheta}.
\]
On the proved box $0\leq h\tan\vartheta<a$, $y<18a$,
$b<4a$, and the denominator of $\mathcal F_3$ exceeds $2/3$.
Thus $|\mathcal F_2|\leq10a^2\sigma_1$ and
$|\mathcal F_3|\leq270a^2\sigma_1+
6aU_{\Omega,2}\vartheta_*$. Their convex combination satisfies
the asserted bound, including the preceding growth interval.

For each layer set the explicit constants
\begin{equation}\label{tts:parentbounds}
 \begin{gathered}
 D_j^*=A_*+\sum_{i<j}U_{\Omega,i},\qquad
 G_j^*=A_0+\sum_{i<j}K_iU_{\Theta,i}r_{+,i},\\
 Q_j^*=\frac{U_{\Omega,j}}{K_j^2r_{-,j}^2},\qquad
 Q_{{\rm r},j}^*=
 \frac{W_{j,1}U_{\Omega,j}+\kappa_jU_{\Theta,j}
                              +2D_j^*U_{\Omega,j}}
      {K_j^2r_{-,j}^2}.
 \end{gathered}
\end{equation}
The rank-one operator
$J\zeta_i\otimes\zeta_i/r_i^2$ has norm one, proving
$\|D_{<j}\|\leq D_j^*$. The gradient sum in
\eqref{tts:material} proves $|G_{<j}|\leq G_j^*$.
The identity $|\dot\rho_j/\rho_j|\leq2D_j^*$ proves
$|Q_j|\leq Q_j^*$ and $|Q_{{\rm r},j}|\leq Q_{{\rm r},j}^*$.
In particular the second-parent velocity contributes the complete
$U_{\Omega,2}$ to $D_3^*$, with no sign-based removal.

Define finite sums
\begin{align}\label{tts:productbounds}
 e_{j,m}&=(N_j/\ell_j)^mf_m,\nonumber\\
 B_{j,m}(W)&=\sum_{q=0}^m\binom mq
                      \kappa_j^q[W]_qe_{j,m-q},\\
 C_{j,m}(P)&=\sum_{q=0}^m\binom mq
                      \kappa_j^q[P]_qe_{j,m-q+1}.\nonumber
\end{align}
The $q$ derivatives allocated to the linear phase have size at
most $\kappa_j^q$; the other $m-q$ derivatives of the linear
material coordinate have size at most $(N_j/\ell_j)^{m-q}$.
The product rule with all $\binom mq$ allocations therefore proves
$|g_j|_m\leq e_{j,m}$,
$|W(\xi_j)g_j|_m\leq B_{j,m}(W)$, and
$|P(\xi_j)\nabla g_j|_m\leq C_{j,m}(P)$.
These bounds hold for every choice of unit derivative directions.

The full damping-plus-diffusion costs are
\begin{align}\label{tts:diffusionnorms}
 \mathfrak D_{\theta,j,m}
 &=dU_{\Theta,j}
       \{\kappa_j^2B_{j,m}(F)+2B_{j,m+2}(F)\},\nonumber\\
 \mathfrak D_{u,j,m}
 &=dQ_j^*
       \{\kappa_j^2B_{j,m+1}(P)+2B_{j,m+3}(P)\}.
\end{align}
They bound, respectively, $-d(\rho_jV_j+\Delta V_j)$ and
$-d(\rho_jv_j+\Delta v_j)$. The explicit factor two bounds
the sum of the two coordinate traces defining the Laplacian;
it is not an identification with a single second derivative.
The complete force estimates are
\begin{align}\label{tts:scalarcost}
 |E_{\theta,j}|_m\leq{}&W_{j,1}U_{\Theta,j}B_{j,m}(F)
  +Q_j^*G_j^*C_{j,m}(P)+\mathfrak D_{\theta,j,m}\nonumber\\
 &+Q_j^*U_{\Theta,j}
    \sum_{q=0}^m\binom mq B_{j,q+1}(P)B_{j,m-q+1}(F),\\
 |E_{u,j}|_m\leq{}&(2D_j^*Q_j^*+Q_{{\rm r},j}^*)B_{j,m+1}(P)
   +U_{\Theta,j}B_{j,m}(F)+\mathfrak D_{u,j,m}\nonumber\\
 &+(Q_j^*)^2\sum_{q=0}^m\binom mq
                    B_{j,q+1}(P)B_{j,m-q+2}(P).
 \label{tts:vectorcost}
\end{align}
For the scalar nonlinear term apply the product rule to
$v_j\cdot\nabla V_j$; for the vector term apply it to
$(\nabla v_j)v_j$. The displayed derivative allocations give
exactly the two sums. All other terms follow individually from
\eqref{tts:scalarforce}--\eqref{tts:spatialjets}. This proves
the estimates with activation, all parent interactions, self-advection,
buoyancy, moving phase length, and diffusion included.

For the base on $I$, let
$c_m=\sup_x\|D^m\chi_0(x)\|$, $H_m=|J\nabla H_0|_m$, and
\[
 B_{0,m}(F)=\sum_{q=0}^m\binom mq\mu^q[F]_qc_{m-q}.
\]
Write $A_{**}$ for the explicit bound on $|\ddot\alpha|$
constructed in \eqref{tts:controljets} below. Then
\begin{align}\label{tts:basecost}
 |f_{\theta,b}|_m&\leq2d(A_0/\mu)B_{0,m+2}(F),\nonumber\\
 |f_{u,b}|_m&\leq A_{**}H_m+
 A_*^2\sum_{q=0}^m\binom mq H_{q+1}H_{m-q}
 +(A_0/\mu)B_{0,m}(F)+2dA_*H_{m+2}.
\end{align}
These follow directly from \eqref{tts:baseforce}, since the
interpolation is complete on $I$. Adding these bounds to
\eqref{tts:scalarcost} and \eqref{tts:vectorcost} gives explicit
bounds for both total forces at every spatial order. The previously
proved formulas cover $[0,t_b]$ with the same initial interpolation;
the maximum of its bounds and these new ones covers $[0,t_1]$.
Every bound on $I$ integrates with the explicit length
$(t_a-t_b)+\Gamma^{-1}<15/\sigma_1$ from \eqref{tts:duration}.

In particular all localization and frequency factors remain visible:
\begin{equation}\label{tts:frequencycost}
 \begin{gathered}
 N_1/\ell_1=C_\ell\mu/s_0,\quad
 N_2/\ell_2=\widetilde N\mu\widehat\lambda_1^3,\quad
 N_3/\ell_3=\frac{16\mu\widehat\lambda_2^3}{\sin s},\\
 dU_{\Theta,3}K_3^2
       =de^4A_0\mu\widehat\lambda_3^{9/8},\qquad
 dQ_3^*K_3^3
       =\frac{64e^4dA_0\mu}{21\sigma_2}
                                      \widehat\lambda_3^{9/8}.
 \end{gathered}
\end{equation}
These are exact factors in upper costs, not lower bounds for the
total force. The exact modal diffusion exponents on $I$ are
\begin{equation}\label{tts:exponents}
 dK_1^2|I|,\qquad dK_2^2\int_I D\,dt,\qquad
 dK_3^2\int_I R\,dt.
\end{equation}
The second lies between $dK_2^2|I|$ and $11dK_2^2|I|$.
The third splits exactly into the proved growth cost and
$\delta_3\int_0^1R\,d\tau$, whose transition part lies between
$(63/64)\delta_3$ and $3\delta_3$, with
$\delta_3=dK_3^2/\Gamma$. Older coupling terms remain in
\eqref{tts:material}; they have not been replaced by damping.

\subsection{An explicit recursion for every temporal and mixed cost}

We give the coefficient recursion as well as the spatial chain rule,
so no stationary-envelope convention enters a time derivative.
Use $\tau=\Gamma(t-t_a)$, $L=L_3$,
$q=\sigma_1/\sigma_2=s/L$, $\epsilon=\sigma_1/\Gamma$,
$\delta_j=dK_j^2/\Gamma$, and the dimensionless state
\begin{equation}\label{tts:state}
 \begin{gathered}
 z=(h,g,u,E,F,\omega,V,X,Y),\\
 E=-K_1\Theta_1/\sigma_1^2,\quad
 F=-K_2\Theta_2/\sigma_2^2,\quad
 \omega=\Omega_1/\sigma_1,\quad V=\Omega_2/\sigma_2,\\
 X=-\Theta_3/P_3^*,\quad Y=-\sigma_2\Omega_3/(K_3P_3^*).
 \end{gathered}
\end{equation}
The two last coordinates are signed linear amplitudes; forming them
does not divide by a temperature. Put
$\bar c=A_0\cos s_*/\sigma_1^2$ and
$\bar C=A_0\sin s_*/\sigma_1^2$, and define
\[
 \begin{gathered}
 x_c=(hg-ab)/D,\quad y_c=(ag+hb)/D,\\
 N_3=\sigma_1^2(E+\bar c)y_c+\sigma_1^2\bar Cx_c
                                      +b\sigma_2^2F,\\
 \quad\mathcal K=\frac{N_3}{\sigma_2^2\sin s\,R},
 \quad\mathcal B=\frac{g\sin(su)+b\cos(su)}{\sqrt D\sin s}.
 \end{gathered}
\]
Let $\Phi(z,\lambda)=(1-\lambda)\mathcal F_2+
\lambda\mathcal F_3$, with the exact feedbacks displayed above
and substitutions $\Omega_1=\sigma_1\omega$,
$\Omega_2=\sigma_2V$, $\vartheta=su$. The complete explicit
dimensionless vector field $\mathcal V(z,\lambda)$ is
\begin{align}\label{tts:vectorfield}
 h'&=a\epsilon\omega,&
 g'&=\epsilon\omega\frac{2ahg+(h^2-a^2)b}{D}
                                  +\frac b{\sin s}V,\nonumber\\
 u'&=\frac{-\Phi-a^2\sigma_1\omega/D}{\Gamma s},&
 E'&=-\epsilon\bar C\omega-\delta_1E,\nonumber\\
 F'&=-q\epsilon\frac{a(E+\bar c)+\bar Ch}{D}V-\delta_2DF,&
 \omega'&=-\epsilon E\frac{h\sin(su)-a\cos(su)}{\sqrt D}
                                               -\delta_1\omega,\\
 V'&=-F\sqrt D\frac{\sin(su)}{\sin s}-\delta_2DV,&
 X'&=\mathcal K Y-\delta_3RX,\nonumber\\
 Y'&=\mathcal B X-\delta_3RY.&&\nonumber
\end{align}
Every prime in this display is $d/d\tau$. Substitution of
\eqref{tts:state} into the physical equations proves each row,
including both factors $q\epsilon$ in the second-temperature row.
Conversely those substitutions reverse, giving the physical pair,
phases, and feedback. This explicitly proves the coordinate map and
its inverse wherever the displayed denominators are positive.

For bounds let $\mathcal Z$ be the closed rectangular box given
by
\[
 \begin{gathered}
 h\in[h_0,13/4],\ g\in[3/4,5],\ u\in[0,1/1024],\\
 E\in[1/8,2],\ F\in[1/16,3],\ \omega\in[0,10],\\
 V\in[-1/32,0],\ X\in[0,e^4],\ Y\in[0,3e^4].
 \end{gathered}
\]
The proved actual solution on $I$ belongs to this box. All
feedback denominators are bounded away from zero on the whole
box by the same elementary inequalities used above. In particular
$D\geq1$ and $g\cos(su)-b\sin(su)>2/3$ there. Introduce
independent symbols $\lambda^{[0]},\lambda^{[1]},\ldots$ for the jets
of the control switch and define on each finite jet expression
\begin{equation}\label{tts:jetderivation}
 \mathcal C=
 \sum_{l=1}^9\mathcal V_l(z,\lambda^{[0]})\partial_{z_l}
             +\sum_{r\geq0}\lambda^{[r+1]}\partial_{\lambda^{[r]}}.
\end{equation}
The second sum is finite on every expression to which it is applied.
The chain rule proves that the physical $r$th time derivative of
any state function $H(z,\lambda)$ is
$\Gamma^r\mathcal C^rH$ evaluated at
$\lambda^{[j]}=\eta^{(j)}(\Gamma(t-t_a))$. Before the transition
these are zero for every $j\geq0$;
the same identity therefore holds on all of $I$.
Consequently the explicit control constants are
\begin{equation}\label{tts:controljets}
 A_{r+1}=\Gamma^r
 \max\left\{|\mathcal C^r\Phi|:
 z\in\mathcal Z,\ 0\leq\lambda^{[0]}\leq1,
 \ |\lambda^{[j]}|\leq[\eta^{(j)}]_0\ (1\leq j\leq r)\right\},
 \qquad A_{**}=A_2.
\end{equation}
This is a maximum of a fully specified continuous elementary
function over an explicit compact rectangle; it is finite and
depends only on the retained parameters and profile derivative
bounds. It is not a supremum over an unspecified later solution.
The equality of derivatives just proved yields
$|\alpha^{(r+1)}|\leq A_{r+1}$ at every order. The sharper
$A_*$ in \eqref{tts:normdata} bounds its first derivative.

For a complete algebraic mixed-jet evaluation, set
$A_j=M_j^{-1}/\ell_j$ and introduce independent
$(t,x,\xi,y)$. Define
\begin{equation}\label{tts:pullbackops}
 \mathfrak T_j=\partial_t+(\dot k_j\cdot x)\partial_\xi
                                  +(\dot A_jx)\cdot\nabla_y,
 \qquad
 \mathfrak X_{j,i}=\partial_{x_i}+k_{j,i}\partial_\xi
                                +\sum_l(A_j)_{li}\partial_{y_l}.
\end{equation}
For $\iota_j^*H=H(t,x,k_j(t)\cdot x,A_j(t)x)$, the ordinary
chain rule, iterated one derivative at a time, gives
\begin{equation}\label{tts:pullbackidentity}
 \partial_t^n\partial_x^\gamma\iota_j^*H
 =\iota_j^*\bigl(\mathfrak T_j^n
        \mathfrak X_{j,1}^{\gamma_1}\mathfrak X_{j,2}^{\gamma_2}H\bigr).
\end{equation}
The commutators vanish: the time derivative of $k_{j,i}$
cancels the $x_i$ derivative of $\dot k_j\cdot x$, and
the two corresponding terms for $A_j$ cancel as well. Thus the
identity also verifies the consistency of all mixed derivative orders.
Every coefficient jet in this recursion is evaluated by
\eqref{tts:vectorfield}--\eqref{tts:controljets} and
\[
 \dot M_j=D_{<j}M_j,\qquad
 \dot A_j=-A_jD_{<j},\qquad
 \dot k_j=-D_{<j}^Tk_j,
\]
together with product differentiation of \eqref{tts:material}
and the exact activation jets $\Gamma^r\eta^{(r)}$ for $w_3$.
For example
$\partial_t^{n+1}A_j=-\sum_{q=0}^n\binom nq
(\partial_t^qA_j)(\partial_t^{n-q}D_{<j})$;
the other two linear equations have the identical product recursion
with their displayed order and signs. Their zeroth bounds are
\eqref{tts:matrixbounds}, so induction gives explicit finite
matrix and phase jet bounds at every prescribed order.

To obtain an actual numerical expression for a mixed norm, expand
\eqref{tts:pullbackidentity} on every displayed field or residual,
keeping its finitely many summands. Replace each coefficient jet by
the preceding recursive bound, each $x^\beta$ by
$(N_j\ell_j)^{|\beta|}$ on its layer support, each periodic
derivative by $[F]_r$ or $[P]_r$, and each envelope derivative by
$f_m$. Sum the absolute products with their explicit binomial
coefficients. The triangle and product inequalities prove the
resulting bound. For the base use its support radius $R_H$ and
its displayed functions instead. This procedure computes all mixed
costs, including every time derivative of $d\rho_j$ and the changing
Laplacian metric; its recursion has finite depth $n+|\gamma|$.

\subsection{The transition endpoint and a proved smooth force collar}

At $t_1$ all weights are one, $\Theta_3(t_1)<0$ and
$\Omega_3(t_1)<0$, with the exact inverse amplitudes of the
transition theorem. In the common affine cell the complete state is
\begin{equation}\label{tts:endpoint}
 \theta=\left(-A_0R_\alpha e_2+
             \sum_{j=1}^3K_j\Theta_j\zeta_j\right)\cdot x,
 \qquad
 u=\left(\dot\alpha J+
    \sum_{j=1}^3\frac{\Omega_j}{r_j^2}
                          J\zeta_j\otimes\zeta_j\right)x.
\end{equation}
The second vorticity has the revived negative sign proved above;
the first remains positive. Formula \eqref{tts:endpoint} and the
full cutoff formulas give actual states of \eqref{tts:PDE}, with
their exact material matrices and laboratory gravity.

To continue smoothly to the right, use the explicit vector field
\eqref{tts:vectorfield} with the constant control value $\lambda=1$.
Let $z_1=z(t_1)$ and take the closed cube
$\mathcal C_1=\{z:\|z-z_1\|_\infty\leq1/100\}$.
Its interior is in the smooth domain of this vector field. Indeed
$h>h_0-1/100>3/4$, $g>3/4-1/100>7/10$, and
$|u|<1/1024+1/100<1/50$ there. Since $s\leq a\leq1/512$
and $b<4s$, these imply $\cos(su)>15/16$ and
\[
 g\cos(su)-b\sin(su)
    >(7/10)(15/16)-4s^2/50>3/5.
\]
Thus $D$, $R$, $\sin s$ and every feedback denominator remain
strictly positive on a neighborhood of the cube. No amplitude
quotient occurs in its vector field. Define the concrete finite maxima
\begin{equation}\label{tts:collarconstants}
 \begin{gathered}
 B=\max_{\mathcal C_1}\|\mathcal V(z,1)\|_\infty,
 \qquad L_{\rm lip}=\max_{\mathcal C_1}
                            \|D_z\mathcal V(z,1)\|_{\infty\to\infty},\\
 \Delta\tau=\min\left\{\frac1{400(1+B)},
                            \frac1{2(1+L_{\rm lip})}\right\}>0.
 \end{gathered}
\end{equation}
On continuous paths in this cube beginning at $z_1$, the Picard
map $z(\tau)\mapsto z_1+\int_1^\tau\mathcal V(z(r),1)dr$
moves points by at most $B\Delta\tau<1/400<1/100$.
The mean-value formula on the convex cube makes its Lipschitz
constant at most $L_{\rm lip}\Delta\tau<1/2$ in the uniform
norm. Successive iterates therefore converge uniformly: their
successive differences are bounded by a geometric series with
ratio below $1/2$. The limit satisfies the integral equation,
and the same contraction inequality proves uniqueness. Differentiation
of this equation first gives $C^1$, and repeated differentiation of
the smooth vector field gives $C^\infty$. This constructs an actual
right collar of physical length
\begin{equation}\label{tts:collarlength}
 \epsilon_R=\Delta\tau/\Gamma>0.
\end{equation}
At $t_1$ the original switching profile is flat and equal to one,
so the control and every derivative match this continuation.
The matrices and phase fields continue by their exact linear ODEs.
Their determinant stays one by the trace identity, or by the
explicit formulas with $\det\mathcal L=-b$. The angle continues
by $\dot\alpha=\mathcal F_3$. Hence the fields have a smooth
right extension with the actual endpoint jets, not a frozen state.

On this collar define the forces by the raw full residuals
\begin{equation}\label{tts:rawcollar}
 f_\theta^{\rm raw}=\theta_t+u\cdot\nabla\theta-d\Delta\theta,
 \qquad
 f_u^{\rm raw}=u_t+(u\cdot\nabla)u-d\Delta u-\theta e_2.
\end{equation}
They agree with \eqref{tts:baseforce} on the entire required
slab by the proof above. Thus no additional nesting assertion
outside that slab is needed for the collar. They are smooth and
compact in space, since the fields are smooth compact fields and
their matrices vary continuously on a compact time interval.
All collar mixed costs are explicitly computed by the same
derivation and pullback recursions, replacing $\mathcal Z$ by
$\mathcal C_1$ and $\lambda$ by one in their maxima, and using
the raw expressions \eqref{tts:rawcollar}. The compact matrix
bounds follow directly from \eqref{tts:matrices} on that cube;
for example $h<13/4+1/100$, $g<5+1/100$, and its positive
denominators bound $\mathcal L$ and $\mathcal L^{-1}$ there.
These explicit finite maxima provide every derivative order without
assuming the original pre-endpoint amplitude bounds on the collar.

Extend the original sine state constantly to negative times, with
zero velocity and inactive increments, and put $\epsilon_L=1/\nu$.
Use the raw residuals of these extensions and set
\begin{equation}\label{tts:timecutoff}
 \beta(t)=\eta(1+t/\epsilon_L)
                   \eta(1+(t_1-t)/\epsilon_R),\qquad
 \widetilde f_\bullet=\beta f_\bullet^{\rm raw}.
\end{equation}
They equal the proved forces on $[0,t_1]$. Flatness of $\eta$
makes their zero extension outside
$[-\epsilon_L,t_1+\epsilon_R]$ smooth. Their support is compact
in space and time. More explicitly, with
\[
 B_{\beta,r}=\sum_{q=0}^r\binom rq
 \epsilon_L^{-q}\epsilon_R^{-(r-q)}
              [\eta^{(q)}]_0[\eta^{(r-q)}]_0,
\]
the product rule gives for every multi-index $\gamma$
\begin{equation}\label{tts:cutoffcost}
 \|\partial_t^n\partial_x^\gamma\widetilde f_\bullet\|_\infty
 \leq\sum_{q=0}^n\binom nq B_{\beta,q}
       \|\partial_t^{n-q}\partial_x^\gamma
                              f_\bullet^{\rm raw}\|_\infty.
\end{equation}
The right-hand mixed norms have the explicit finite recursion just
proved. The PDE is asserted on $[0,t_1]$, where the cutoff is one.

Finally, the streamfunction is compact and smooth; hence the velocity
has finite energy. If $N_0(x)=(2\pi)^{-1}\log|x|$, then
$N_0*\Delta\psi=(\Delta N_0)*\psi=\psi$ by distributional
integration by parts, and therefore
$u=J\nabla N_0*\omega$, $\omega=\Delta\psi$, exactly.
Odd $F$, even $P$, and even material envelopes make
$\theta,u,f_\theta,f_u$ odd and $\omega$ even. Every field and
force on the required slab has the original spatial support inside
$B_{R_H}$ by the proved nesting. This completes the full compact
spatial extension through the actual third transition with every
retained physical diffusion cost and both older interaction directions.
